\chapter{Price and Volume Features}\label{ch:rs:price-and-volume-features}

The momentum factor published in the Kenneth French data library, like most momentum signals, skips the most recent month: it
sorts stocks on their return from twelve months ago to one month ago. The reason is an effect, not a convention. Stocks' monthly returns have a strongly negative
first-order autocorrelation (Jegadeesh, 1990), so the last month's winners tend to give some of it back, and including that
month mixes a reversal into a continuation signal. On the book's synthetic market the skip does the opposite: momentum over
the full twelve months has a one-month \omterm{def:rs:anatomy-of-a-predictor:ic}{information coefficient} of 0.034, the skipped version 0.030. The simulator's reversal
lasts one day, not one month, and in a market without a monthly reversal the skip throws information away. A feature's
details encode claims about the market; this chapter builds the standard family of price and volume features, measures each
on the synthetic market, and says which claim each one makes.

\section{Returns at many scales}

\begin{definition}[Return feature, skip period]\label{def:rs:price-and-volume-features:return}
A \emph{return feature}\index{return feature} is a security's compounded return over a window ending at or before the decision
time, used as a \omterm{def:rs:anatomy-of-a-predictor:predictor}{predictor}. Its \emph{skip period}\index{skip period} is the gap between the window's end and the decision time:
a 12-1 momentum feature uses the return from 252 to 21 trading days ago, skipping the last 21.
\end{definition}

\omterm{def:rs:price-and-volume-features:return}{Return features} are the oldest \omterm{def:rs:anatomy-of-a-predictor:predictor}{predictors} and still the most used: a short window (a day to a month) for reversal, a long one
(three to twelve months) for momentum, three to five years for long-term reversal (De Bondt and Thaler, 1985). Jegadeesh and Titman
(1993) found that buying past winners and selling past losers over three to twelve months earned significant returns; Book~8
studies the strategies. Here the question is the measurement.

The table measures the chapter's twelve features on the synthetic market: the mean rank IC against the next day, week and month,
each sampled once per horizon so that the targets do not overlap and the naive $t$-statistics are valid. The simulator plants a
one-day reversal and a persistent drift (chapter~5), and the table finds both: reversal features score at short horizons, momentum
and trend features at long ones. It plants nothing on volatility, volume or liquidity, and those features score nothing
distinguishable from zero. That is not a finding about real markets, where each has a literature; it is a check that the
measurement does not invent effects.

\begin{center}
\small
\begin{tabular}{lrrr}
\toprule
feature (known at the close of $t$) & 1 day & 5 days & 21 days \\
\midrule
reversal, 1 day & 0.0378 (16.2) & 0.0150 (3.0) & $-0.0058$ ($-0.6$) \\
reversal, 5 days & 0.0176 (8.1) & 0.0085 (1.8) & $-0.0162$ ($-1.8$) \\
reversal, 21 days & 0.0022 (1.0) & $-0.0117$ ($-2.6$) & $-0.0285$ ($-2.9$) \\
momentum 12-1 & 0.0076 (1.4) & 0.0167 (1.4) & 0.0296 (1.2) \\
momentum 12 & 0.0059 (1.1) & 0.0177 (1.5) & 0.0340 (1.4) \\
momentum 6-1 & 0.0067 (1.8) & 0.0140 (1.7) & 0.0285 (1.6) \\
low volatility (21 days) & $-0.0034$ ($-2.5$) & $-0.0033$ ($-1.1$) & $-0.0023$ ($-0.5$) \\
\omterm{def:rs:price-and-volume-features:volume}{volume surprise} & 0.0009 (1.3) & $-0.0007$ ($-0.5$) & $-0.0008$ ($-0.3$) \\
\omterm{def:rs:price-and-volume-features:volume}{Amihud illiquidity} & $-0.0004$ ($-0.3$) & $-0.0023$ ($-0.7$) & $-0.0055$ ($-0.8$) \\
low turnover & 0.0011 (1.5) & 0.0026 (1.7) & 0.0049 (1.7) \\
crossover 5/20 & 0.0027 (1.2) & 0.0133 (2.9) & 0.0218 (2.4) \\
crossover 50/200 & 0.0084 (2.0) & 0.0173 (1.9) & 0.0324 (1.7) \\
\bottomrule
\end{tabular}
\end{center}
\noindent{\small Mean rank IC (and $t$-statistic) of each feature against the next 1, 5 and 21 days' returns on the synthetic market,
ten years, one observation per horizon. Data: \texttt{firm.synthmkt}, seed 1.}

\begin{proposition}[What the skip is worth]\label{prop:rs:price-and-volume-features:skip}
Write the twelve-month log return as the sum of the 12-1 part and the last month's. If both are standardised the same way, the
IC of the full return is, to first order, a variance-weighted average of the two parts' ICs. The skip raises the IC exactly when the
last month's return has a lower IC than the 12-1 part, in particular when it is negative (a monthly reversal).
\end{proposition}

\begin{proof}
For a sum $s = a + b$ and a target $y$, $\operatorname{corr}(s, y) = (\sigma_a\operatorname{corr}(a, y) + \sigma_b\operatorname{corr}(b,
y))/\sigma_s$, with $\sigma_s \le \sigma_a + \sigma_b$; the full signal beats $a$ alone if and only if $\operatorname{corr}(b, y)$ is large
enough relative to $\operatorname{corr}(a, y)$, and it loses whenever $\operatorname{corr}(b, y) < 0$.
\end{proof}

On the synthetic market the last month's own return has an IC of $+0.0285$ with the next month (the reversal feature's $-0.0285$ with
the sign changed): continuation, so the full year wins. In the US data of Jegadeesh (1990) the last month reverses, and the skip wins.
\Cref{fig:rs:price-and-volume-features:horizon} shows each feature's IC against the length of its target: the one-day reversal is
spent within a week, momentum keeps accumulating for three months.

\begin{omfigure}
\begin{tikzpicture}
\begin{semilogxaxis}[width=10.5cm,height=5.2cm,xlabel={target horizon $h$ (trading days, log scale)},ylabel={mean rank IC},
  tick label style={font=\small},label style={font=\small},xmin=0.9,xmax=70,ymin=-0.03,ymax=0.045,grid=major,grid style={gray!25},
  xtick={1,2,5,10,21,42,63},xticklabels={1,2,5,10,21,42,63},scaled y ticks=false,
  yticklabel style={/pgf/number format/fixed,/pgf/number format/precision=2},legend columns=2,legend cell align=left,
  legend style={font=\footnotesize,at={(0.5,-0.32)},anchor=north,draw=none,fill=none,
  /tikz/every even column/.append style={column sep=8pt}},table/col sep=comma]
\addplot[omThm,very thick,mark=*] table[x=h,y=rev1]{figdata/research/07-price-and-volume-features/horizon.csv};
\addplot[omThm!55,thick,dashed,mark=o] table[x=h,y=rev21]{figdata/research/07-price-and-volume-features/horizon.csv};
\addplot[omDef,very thick,mark=square*] table[x=h,y=mom121]{figdata/research/07-price-and-volume-features/horizon.csv};
\addplot[omDef!55,thick,dashed,mark=square] table[x=h,y=mom12]{figdata/research/07-price-and-volume-features/horizon.csv};
\addplot[black!60,thin,sharp plot] coordinates {(0.9,0) (70,0)};
\legend{reversal 1 day,reversal 21 days,momentum 12-1,momentum 12}
\end{semilogxaxis}
\end{tikzpicture}

\omcaption{Mean rank IC of four \omterm{def:rs:price-and-volume-features:return}{return features} against returns over the next $h$ days on the synthetic market, scored every day.
The one-day reversal peaks at a day and is gone after two weeks; the last month's reversal is negative at every horizon (the month
continues); momentum rises to about 0.04 at a quarter, slightly higher without the skip. Data: \texttt{firm.synthmkt}, seed 1.}\label{fig:rs:price-and-volume-features:horizon}
\end{omfigure}

\section{Volatility and range}

\begin{definition}[Range-based volatility estimator]\label{def:rs:price-and-volume-features:range}
A \emph{range-based volatility estimator}\index{range-based volatility estimator} estimates a day's return variance from its open,
high, low and close prices rather than from close-to-close returns alone.
\end{definition}

\begin{definition}[Parkinson, Garman--Klass and Yang--Zhang estimators]\label{def:rs:price-and-volume-features:estimators}
With $H, L, O, C$ a day's high, low, open and close, the \emph{Parkinson estimator}\index{Parkinson estimator} of the daily variance is
$(\ln H/L)^2/(4\ln 2)$; the \emph{Garman--Klass estimator}\index{Garman--Klass estimator} is $\tfrac12(\ln H/L)^2 - (2\ln 2 -
1)(\ln C/O)^2$; the Rogers--Satchell estimator is $\ln(H/C)\ln(H/O) + \ln(L/C)\ln(L/O)$. The \emph{Yang--Zhang
estimator}\index{Yang--Zhang estimator} over $n$ days adds the sample variance of overnight returns, $k$ times that of open-to-close
returns and $1 - k$ times the average Rogers--Satchell term, with $k = 0.34/(1.34 + (n+1)/(n-1))$.
\end{definition}

For a Brownian motion without drift observed continuously, $\E[(\ln H/L)^2] = 4\ln2\,\sigma^2$, which is Parkinson's normalisation
(Parkinson, 1980); the range carries much more information about $\sigma$ than the close alone. Each later estimator repairs a
weakness of the one before. Garman and Klass (1980) combine the range with the open-to-close return; Rogers and Satchell (1991) are
unbiased whatever the drift; Yang and Zhang (2000) add the overnight gap, which no intraday range can see.

The chapter's simulation of 4\,000 days of one-minute Brownian paths scores them (\cref{fig:rs:price-and-volume-features:race}).
Without drift or gaps, the variance of the \omterm{def:rs:price-and-volume-features:estimators}{Parkinson estimator} is 5.2 times smaller than that of the squared close-to-close return,
Garman--Klass's 7.6 times, Rogers--Satchell's 5.7 times; Yang--Zhang over twenty days is 7.5 times more efficient than the
twenty-day close-to-close variance. Every range estimator is about 8--10\% low, because a range observed at 390 points is narrower
than the continuous one: discrete monitoring is a bias the formulas do not know about. A drift of half a standard deviation per day
inflates the squared close-to-close return by 25\% and Parkinson by 2\%, leaves Rogers--Satchell and Yang--Zhang unchanged; an overnight
gap carrying 30\% of the variance makes the three single-day range estimators see only 63--64\% of it, and only Yang--Zhang recovers
the total.

\begin{omfigure}
\begin{tikzpicture}
\begin{axis}[width=11cm,height=5.4cm,ybar,bar width=8pt,ylabel={mean estimate / true variance},tick label style={font=\small},
  label style={font=\small},xmin=-0.6,xmax=4.6,ymin=0,ymax=1.35,xtick={0,1,2,3,4},
  xticklabels={{close-to-\\close},{Parkinson},{Garman--\\Klass},{Rogers--\\Satchell},{Yang--\\Zhang}},
  xticklabel style={align=center,font=\footnotesize},
  grid=major,grid style={gray!25},legend columns=3,legend style={font=\footnotesize,at={(0.5,-0.36)},anchor=north,draw=none,
  fill=none,/tikz/every even column/.append style={column sep=8pt}},table/col sep=comma]
\addplot[fill=omDef!55,draw=omDef] table[x=k,y=base]{figdata/research/07-price-and-volume-features/race.csv};
\addplot[fill=omThm!45,draw=omThm] table[x=k,y=drift]{figdata/research/07-price-and-volume-features/race.csv};
\addplot[fill=black!25,draw=black!60] table[x=k,y=gaps]{figdata/research/07-price-and-volume-features/race.csv};
\addplot[black!70,thin,sharp plot,forget plot] coordinates {(-0.6,1) (4.6,1)};
\legend{no drift or gap,drift of half a daily sd,overnight gap: 30\% of variance}
\end{axis}
\end{tikzpicture}

\omcaption{Mean of each daily variance estimator relative to the true diffusion variance, on 4\,000 simulated days of 390 one-minute
steps. All range estimators read 8--10\% low from discrete monitoring; a drift inflates close-to-close and Parkinson; an overnight gap
is invisible to the single-day range estimators and recovered by Yang--Zhang. Data: the chapter's module, seeded.}\label{fig:rs:price-and-volume-features:race}
\end{omfigure}

\section{Volume and liquidity}

\begin{definition}[Volume surprise, Amihud illiquidity, turnover ratio]\label{def:rs:price-and-volume-features:volume}
The \emph{volume surprise}\index{volume surprise} of a security on a day is the logarithm of its volume divided by its average volume
over a preceding window. \emph{Amihud illiquidity}\index{Amihud illiquidity} is the average over a window of the absolute daily return
divided by the day's traded value: the price move per unit of money traded (Amihud, 2002). The \emph{turnover ratio}\index{turnover
ratio} is traded volume divided by shares outstanding, averaged over a window.
\end{definition}

Volume features measure attention and liquidity, and each has a claim behind it. Gervais, Kaniel and Mingelgrin (2001) found that
stocks with unusually high volume over a day or a week tend to rise over the following month (the high-volume return premium): volume
draws attention and investors' attention moves prices. Amihud's measure is a cheap proxy for price impact, and illiquid stocks are
expected to pay a premium for bearing it. None of these mechanisms is in the book's simulator, whose volume follows the size of each
day's move and the volatility state but carries no information about future returns; the table shows exactly that. A researcher who
wants to study them needs real data, or a simulator that plants them.

\section{Technical families as linear filters}

\begin{definition}[Moving-average crossover, linear filter]\label{def:rs:price-and-volume-features:filter}
A \emph{moving-average crossover}\index{moving-average crossover} with windows $n < m$ is the logarithm of the ratio of the average of
the last $n$ prices to the average of the last $m$ prices; the classic rule is long when it is positive. A \emph{linear
filter}\index{linear filter} of returns is a \omterm{def:rs:anatomy-of-a-predictor:predictor}{predictor} of the form $\sum_{j \ge 0} w_jr_{t-j}$ with fixed weights $w_j$.
\end{definition}

\begin{proposition}[A crossover is a filter of past returns]\label{prop:rs:price-and-volume-features:crossover}
To first order in the returns, $\ln(\overline P_n/\overline P_m) = \sum_{j=0}^{m-2} w_jr_{t-j}$ with $w_j = \frac{(m - 1 - j)}{m} -
\frac{(n - 1 - j)^+}{n}$, where $r_{t-j}$ is the log return of $j$ days ago: the weights rise linearly over the short window, then fall
linearly to zero at the long window's length, and sum to $\frac{m - n}{2}$.
\end{proposition}

\begin{proof}
The log of an average of prices is, to first order, the average of the log prices. The average of the last $n$ log prices is $\ln P_t
- \sum_{j<n-1}\frac{n-1-j}{n}r_{t-j}$, since the price $k$ days ago is $\ln P_t$ minus the $k$ most recent returns. Subtract the
same expression with $m$. The weights sum to $\frac{m-1}{2} - \frac{n-1}{2}$.
\end{proof}

A crossover is therefore a momentum signal with a triangular window that skips most of the last $n$ days
(\cref{fig:rs:price-and-volume-features:filter}), and its apparent variety is a variety of weights on past returns. Brock, Lakonishok
and LeBaron (1992) tested moving-average and trading-range rules on the Dow Jones index from 1897 to 1986 and found strong support for
them; the filter view says what such a finding is about: the autocorrelation of index returns at the lags the weights emphasise. On the
synthetic market, the 5/20 crossover has an IC of 0.022 at a month ($t = 2.4$), a smoothed version of the momentum the simulator plants.

\begin{omfigure}
\begin{tikzpicture}
\begin{axis}[width=10.5cm,height=4.8cm,xlabel={lag $j$ of the daily return (days)},ylabel={weight $w_j$},tick label style={font=\small},
  label style={font=\small},xmin=0,xmax=50,ymin=0,ymax=0.9,grid=major,grid style={gray!25},legend columns=2,
  legend style={font=\footnotesize,at={(0.5,-0.32)},anchor=north,draw=none,fill=none,
  /tikz/every even column/.append style={column sep=8pt}},table/col sep=comma]
\addplot[omDef,very thick,const plot mark mid] table[x=lag,y=w520]{figdata/research/07-price-and-volume-features/filter.csv};
\addplot[omThm,very thick,dashed,const plot mark mid] table[x=lag,y=w1050]{figdata/research/07-price-and-volume-features/filter.csv};
\legend{crossover 5/20,crossover 10/50}
\end{axis}
\end{tikzpicture}

\omcaption{Weights a \omterm{def:rs:price-and-volume-features:filter}{moving-average crossover} puts on each past daily return (\cref{prop:rs:price-and-volume-features:crossover}):
rising over the short window, then falling linearly to zero at the long window. The 5/20 crossover sums to 7.5 and peaks at lag 5; the
10/50 to 20 and peaks at lag 10.}\label{fig:rs:price-and-volume-features:filter}
\end{omfigure}

\section{Predictor cards}

\begin{predictorcard}{Momentum 12-1}\label{pred:rs:price-and-volume-features:momentum}
\sfield{Definition} $s_{i,t} = \prod_{k=t-251}^{t-21}(1 + r_{i,k}) - 1$: the return from 252 to 21 trading days ago, ranked.
\sfield{Inputs and timestamps} Daily closes up to 21 days before the decision date; \omterm{def:rs:universe-symbology-and-corporate-actions:delisting}{delisting returns} included.
\sfield{Rationale} Underreaction to information and slow diffusion; the last month is skipped because it reverses in US data
(Jegadeesh, 1990).
\sfield{Horizon and half-life} One to three months. On the synthetic market the IC with a single day's return is 0.0076 the next day and
about 0.004 six months later; an exponential fit gives a half-life of 260 days.
\sfield{Normalisation} Cross-sectional rank; neutralised to industries when industry momentum is not wanted.
\sfield{Failure modes} Momentum crashes after market rebounds (Book~8); crowding; turnover costs.
\sfield{Sources} Jegadeesh and Titman (1993); \texttt{rs\_features.lag\_profile} and \texttt{half\_life}.
\end{predictorcard}

\begin{predictorcard}{Volume surprise}\label{pred:rs:price-and-volume-features:volume}
\sfield{Definition} $s_{i,t} = \ln\bigl(V_{i,t}/\tfrac1{21}\sum_{k=1}^{21}V_{i,t-k}\bigr)$, ranked.
\sfield{Inputs and timestamps} Daily volumes to the close of $t$ (consolidated volume is known after the close).
\sfield{Rationale} Unusual volume attracts attention and buyers; the high-volume return premium (Gervais, Kaniel and Mingelgrin, 2001).
\sfield{Horizon and half-life} A month in the literature. Not measurable on the synthetic market, which plants no volume effect: its IC
is 0.0009 ($t = 1.3$) at a day and zero after.
\sfield{Normalisation} Rank; the average excludes the current day.
\sfield{Failure modes} Volume on news carries the news' direction, not attention; index-rebalance and option-expiry days inflate volume
mechanically.
\sfield{Sources} Gervais, Kaniel and Mingelgrin (2001); the chapter's table.
\end{predictorcard}

\section{Tutorial: the feature library}\label{tut:rs:price-and-volume-features:library}

\begin{tutorial}
\textbf{Goal.} Compute the chapter's twelve features on the synthetic market, prove that none reads the future, and fill the IC table.
\textbf{End state:} the table; \cref{fig:rs:price-and-volume-features:horizon,fig:rs:price-and-volume-features:race,fig:rs:price-and-volume-features:filter}.
\begin{tutsteps}
\item \textbf{Windows without look-ahead.} Every rolling statistic is built from cumulative sums: row $t$ uses rows $t - w + 1$ to $t$,
and a window with a missing value is missing.
\omcode{code/firm/features/firm_features.py}{29}{53}{Rolling sums, shifts and past returns.}\label{lst:rs:price-and-volume-features:rolling}
\item \textbf{The leakage test.} Perturb every input after a cut-off and require the feature up to the cut-off to be unchanged; the
acceptance tests run it on every feature, and on a feature that peeks, which it must catch.
\omcode{code/firm/features/firm_features.py}{137}{148}{The leakage test.}\label{lst:rs:price-and-volume-features:leak}
\item \textbf{Run} \texttt{ic\_table()}, \texttt{ic\_by\_horizon} for the four \omterm{def:rs:price-and-volume-features:return}{return features}, \texttt{range\_race} for the three cases
and \texttt{fig\_features.py}.
\end{tutsteps}
\textbf{What to change next.} Add long-term reversal (the return from five years to one year ago) and check its IC; replace the
simulated minute paths by a firm.tape day and compare the range estimators on a price with a bid--ask bounce.
\end{tutorial}

\section{Build: the feature library}\label{bld:rs:price-and-volume-features:features}

\begin{build}
\bfield{Purpose} The miniature firm's standard \omterm{def:rs:market-data-for-research:bar}{bar} features, computed one way everywhere, with a test that none of them looks ahead.
\bfield{Interface} \texttt{past\_return(ret, window, skip)}, \texttt{rolling\_vol}, \texttt{ewma\_vol}, \texttt{parkinson},
\texttt{garman\_klass}, \texttt{rogers\_satchell}, \texttt{yang\_zhang}, \texttt{volume\_surprise}, \texttt{amihud},
\texttt{turnover}, \texttt{ma\_crossover}, \texttt{crossover\_weights}, \texttt{leakage\_test(feature, inputs, cut)}.
\bfield{Rules} Panels (dates, names); row $t$ of an output uses input rows up to $t$ only; windows with a missing value are missing;
volume averages exclude the current day.
\bfield{Acceptance tests} \texttt{code/firm/features/tests/}: past returns with a skip; every feature passes the leakage test and a
peeking feature fails it; range estimators within 10\% of the true variance of a simulated Brownian day; the crossover equals its
filter to first order.
\bfield{Stretch} Intraday features from firm.tape \omterm{def:rs:market-data-for-research:bar}{bars}; features on volume-clock \omterm{def:rs:market-data-for-research:bar}{bars} (chapter~2); a registry that records each
feature's window and \omterm{def:rs:point-in-time-data-and-the-biases:bitemporal}{knowledge time} for the \omterm{def:rs:anatomy-of-a-predictor:card}{predictor card}.
\end{build}

\begin{omsources}
\item N.~Jegadeesh, ``Evidence of predictable behavior of security returns'', \emph{Journal of Finance} 45(3), 1990; N.~Jegadeesh and
S.~Titman, ``Returns to buying winners and selling losers'', \emph{Journal of Finance} 48(1), 1993.
\item W.~F.~M.~De Bondt and R.~Thaler, ``Does the stock market overreact?'', \emph{Journal of Finance} 40(3), 1985.
\item M.~Parkinson, ``The extreme value method for estimating the variance of the rate of return'', \emph{Journal of Business} 53(1),
1980; M.~B.~Garman and M.~J.~Klass, ``On the estimation of security price volatilities from historical data'', \emph{Journal of
Business} 53(1), 1980.
\item L.~C.~G.~Rogers and S.~E.~Satchell, ``Estimating variance from high, low and closing prices'', \emph{Annals of Applied
Probability} 1(4), 1991; D.~Yang and Q.~Zhang, ``Drift-independent volatility estimation based on high, low, open, and close prices'',
\emph{Journal of Business} 73(3), 2000.
\item Y.~Amihud, ``Illiquidity and stock returns: cross-section and time-series effects'', \emph{Journal of Financial Markets} 5(1),
2002; S.~Gervais, R.~Kaniel and D.~H.~Mingelgrin, ``The high-volume return premium'', \emph{Journal of Finance} 56(3), 2001.
\item W.~Brock, J.~Lakonishok and B.~LeBaron, ``Simple technical trading rules and the stochastic properties of stock returns'',
\emph{Journal of Finance} 47(5), 1992.
\end{omsources}

\section{Exercises}

\begin{exercise}[$\star$]\label{exo:rs:price-and-volume-features:1}
A stock's day has $H = 51.20$, $L = 49.80$. What daily volatility does the \omterm{def:rs:price-and-volume-features:estimators}{Parkinson estimator} give?
\end{exercise}

\begin{exercise}[$\star$]\label{exo:rs:price-and-volume-features:2}
Monthly returns of 2\%, $-1\%$, 4\% and 3\% (oldest first). What is the 3-1 momentum feature at the end of the fourth month, and the
four-month return?
\end{exercise}

\begin{exercise}[$\star$]\label{exo:rs:price-and-volume-features:3}
A stock trades USD~40 million a day and moves 1.2\% on average in absolute value. What is its \omterm{def:rs:price-and-volume-features:volume}{Amihud illiquidity}, in percent per
USD~million? And a stock trading USD~2 million that moves 2.5\%?
\end{exercise}

\begin{exercise}[$\star\star$]\label{exo:rs:price-and-volume-features:4}
Compute the weights of the 3/6 crossover by \cref{prop:rs:price-and-volume-features:crossover}, and their sum.
\end{exercise}

\begin{exercise}[$\star\star$]\label{exo:rs:price-and-volume-features:5}
The \omterm{def:rs:price-and-volume-features:estimators}{Parkinson estimator} is 5.2 times as efficient as a squared daily return. How many days of Parkinson estimates give a daily-variance
estimate as precise as sixty days of squared close-to-close returns?
\end{exercise}

\begin{exercise}[$\star\star$]\label{exo:rs:price-and-volume-features:6}
By \cref{prop:rs:price-and-volume-features:skip}, with standard deviations of the 12-1 part and of the last month's return in the ratio
$\sqrt{11}:1$, ICs of 0.03 for the first and $-0.03$ for the second, and the two parts uncorrelated, what is the IC of the full
twelve-month return?
\end{exercise}

\begin{exercise}[$\star\star\star$]\label{exo:rs:price-and-volume-features:7}
\emph{Coding.} Rerun \texttt{range\_race} with 30 steps a day instead of 390. How large does the discrete-monitoring bias of Parkinson
become, and why?
\end{exercise}

\begin{exercise}[$\star\star\star$]\label{exo:rs:price-and-volume-features:8}
\emph{Find the flaw.} ``Our volume-surprise \omterm{def:rs:anatomy-of-a-predictor:predictor}{predictor} divides today's volume by the average volume of the 21 days centred on today, so
that it is not biased by trends in volume.''
\end{exercise}

\section{Problem: The Range-Estimator Race}

\begin{problem}[{Weekend problem --- which daily variance estimator to trust, and when}]\label{pb:rs:price-and-volume-features:1}
The chapter's simulation: 4\,000 days, 390 one-minute Brownian steps a day, a daily volatility of 1\%, seed 5; in turn no drift and no
gap, a drift of half a daily standard deviation, and an overnight gap carrying 30\% of the daily variance.

\medskip\noindent\textbf{Part I --- The formulas.}
\begin{enumerate}
\item Why is Parkinson's constant $4\ln2$?
\item Which estimator needs only the high and the low?
\item Which is unbiased whatever the drift, and why?
\item What does Yang--Zhang add, and what does its weight $k$ equal for twenty days?
\end{enumerate}

\medskip\noindent\textbf{Part II --- Efficiency.}
\begin{enumerate}[resume]
\item What are the efficiencies of Parkinson, Garman--Klass and Rogers--Satchell relative to the squared close-to-close return?
\item And of twenty-day Yang--Zhang relative to the twenty-day close-to-close variance?
\item How many days of close-to-close returns does one day of Garman--Klass replace?
\end{enumerate}

\medskip\noindent\textbf{Part III --- Bias.}
\begin{enumerate}[resume]
\item What are the four range estimators' biases without drift or gap, and what causes them?
\item What does the drift do to close-to-close, to Parkinson and to Rogers--Satchell?
\item What does the gap do to the three single-day range estimators, and to Yang--Zhang?
\item Why does the close-to-close estimator suffer from the drift but not from the gap?
\item How would a bid--ask bounce in the high and low affect the range estimators?
\end{enumerate}

\medskip\noindent\textbf{Part IV --- Choice.}
\begin{enumerate}[resume]
\item For a stock with large overnight gaps, which estimator?
\item For a 24-hour market with no gap and strong intraday trends, which?
\item For a volatility feature that must be computed from daily closes only, which, and at what cost?
\item What would you check before using a range estimator on real data?
\item Why do the efficiency gains matter for a volatility \omterm{def:rs:anatomy-of-a-predictor:predictor}{predictor}?
\item State the \emph{named result}: the efficiencies of the four range estimators and the share of variance the single-day ones miss
with a 30\% gap.
\item Which of the chapter's features use a volatility estimate?
\item In one sentence: what does the range know that the close does not?
\end{enumerate}
\end{problem}

\section{Interview questions}

\begin{interviewq}[$\star$ \iqroles{researcher, trader}]\label{iq:rs:price-and-volume-features:1}
Why do momentum signals skip the most recent month?
\end{interviewq}

\begin{interviewq}[$\star\star$ \iqroles{researcher, risk}]\label{iq:rs:price-and-volume-features:2}
How would you estimate a stock's daily volatility from one month of daily \omterm{def:rs:market-data-for-research:bar}{bars} as precisely as possible?
\end{interviewq}

\begin{interviewq}[$\star\star$ \iqroles{researcher}]\label{iq:rs:price-and-volume-features:3}
Show that a \omterm{def:rs:price-and-volume-features:filter}{moving-average crossover} is a weighted sum of past returns. What weights?
\end{interviewq}

\begin{interviewq}[$\star\star$ \iqroles{developer, mle}]\label{iq:rs:price-and-volume-features:4}
How do you test automatically that a feature pipeline never uses future data?
\end{interviewq}

\begin{interviewq}[$\star\star$ \iqroles{researcher, trader}]\label{iq:rs:price-and-volume-features:5}
What does Amihud's illiquidity measure capture, and what are its weaknesses?
\end{interviewq}

\begin{interviewq}[$\star\star\star$ \iqroles{researcher}]\label{iq:rs:price-and-volume-features:6}
A new feature has an IC of zero on your simulator and 0.02 on real data. What do you conclude about the feature, and about the simulator?
\end{interviewq}
